\theory{Theory of equations}{When can polynomial equations be solved exactly, and when is a formula using radicals impossible?}{The classical theory of equations studies polynomial equations, their roots, formulas, and the algebraic structures that control solvability. Its central historical problem was solved by Galois theory, which replaced formulas with symmetry groups.}{Algebra, number theory, algebraic geometry, numerical root finding, physics, and the study of polynomial systems.}{Roots, multiplicities, resultants, and Galois symmetries explain both exact solvability and the limits of formulas.}

\paragraph{The historical problem}
For a polynomial equation, mathematicians first sought expressions for roots using addition, subtraction, multiplication, division, and radicals. Linear, quadratic, cubic, and quartic equations received formulas in antiquity and the Renaissance. Scipione del Ferro and Tartaglia found cubic methods; Cardano published them, and Ferrari supplied the quartic solution. Bombelli explained the complex numbers appearing in cubic formulas \cite{equations_ref1}.

The general quintic resisted radicals. Ruffini gave an incomplete impossibility argument, and Niels Abel proved in 1824 that a general fifth-degree equation cannot be solved by radicals. Évariste Galois then gave a criterion: the equation is solvable by radicals exactly when its Galois group is solvable \cite{equations_ref2}.

\end{historylist}
\section*{3. Theory}
\subsection*{Core theory}
\paragraph{Roots and systems}
The fundamental theorem of algebra says every nonconstant complex polynomial has a complex root, counted with multiplicity. Vieta's formulas relate coefficients to symmetric sums and products of roots. The discriminant detects repeated roots: it is zero precisely when the polynomial and its derivative share a root.

Linear systems were solved in antiquity; Cramer gave a general determinant formula in the eighteenth century. Today Gaussian elimination is usually more efficient. Systems of polynomial equations lead to resultants, Grobner bases, and algebraic geometry. Integer solutions are studied as Diophantine equations in number theory \cite{equations_ref3}.

\paragraph{A worked example}
For
\[
x^2-5x+6=0,
\]
factoring gives $(x-2)(x-3)=0$, so the roots are $2$ and $3$. Vieta's
formulas explain the calculation without factoring: the roots must add to
$5$ and multiply to $6$. The discriminant is
\[
(-5)^2-4(1)(6)=1,
\]
which is positive, so the two real roots are distinct.

The same questions become harder for a cubic or quintic. A numerical
algorithm may approximate the roots, but the theory of equations asks a
different question: is there a finite expression using radicals at all?
Galois theory answers that question by studying how the roots can be
permuted while leaving the coefficients fixed. Thus the elementary quadratic
example is a model of the larger subject: coefficients constrain roots,
symmetry explains the constraint, and the discriminant detects a change in
the root pattern \cite{equations_ref1,equations_ref2}.

\section*{4. Important theorems and results}
\begin{enumerate}[leftmargin=*,itemsep=0.35em]
\item \textbf{Fundamental theorem of algebra.} Every nonconstant complex polynomial has a complex root and factors completely into linear factors.

\item \textbf{Abel--Ruffini theorem.} The general polynomial of degree five or greater is not solvable by radicals.

\item \textbf{Galois criterion.} A polynomial is solvable by radicals exactly when its Galois group is solvable.

\item \textbf{Vieta's formulas.} Coefficients determine elementary symmetric functions of the roots.

\item \textbf{Discriminant criterion.} A polynomial has a repeated root exactly when its discriminant is zero.

\item \textbf{Bezout's theorem.} Projective plane curves of degrees $m$ and $n$ meet in $mn$ points counted with multiplicity under suitable hypotheses.
\end{enumerate}
\noindent\emph{Source for these results:} \cite{equations_ref1}.

\theoryapplications
\theoryconnections{theory_of_equations}{%
\item Algebra supplies operations on coefficients, complex numbers guarantee that a polynomial factors into linear factors after passing to a suitable field, and linear algebra handles systems and multiplicities.
\item Symmetry enters because permuting the roots can preserve every coefficient; the set of such permutations is the Galois group.
\item Polynomial geometry then interprets several equations as a shape whose intersections and singularities encode the solutions \cite{equations_ref1,equations_ref2}.
}{%
\item The theory feeds Galois theory by turning the question “can roots be written by radicals?” into a question about solvable groups.
\item It feeds algebraic geometry through varieties and elimination, number theory through integer and rational roots, and numerical analysis through algorithms that approximate roots when no useful formula exists.
\item In physics, characteristic equations determine allowed frequencies or energy levels, so root structure becomes a prediction about a system \cite{equations_ref2,equations_ref3}.
}
\section*{Glossary}
\begin{description}[leftmargin=*,style=nextline,itemsep=0.3em]
\item[\textbf{Polynomial equation.}] An equation of the form $a_nx^n + \cdots + a_1x + a_0 = 0$ where the unknown appears only with non-negative integer powers.
\item[\textbf{Root.}] A value of $x$ that satisfies a polynomial equation; also called a zero of the polynomial.
\item[\textbf{Multiplicity.}] The number of times a root appears as a factor in the complete factorization of a polynomial.
\item[\textbf{Discriminant.}] A quantity computed from the coefficients of a polynomial that is zero exactly when the polynomial has a repeated root.
\item[\textbf{Galois group.}] The group of permutations of a polynomial's roots that leave all coefficients fixed; it controls whether the polynomial is solvable by radicals.
\item[\textbf{Solvable by radicals.}] A property of a polynomial equation meaning its roots can be expressed using a finite combination of the coefficients and the operations of addition, subtraction, multiplication, division, and extraction of roots.
\item[\textbf{Vieta's formulas.}] Relations expressing the coefficients of a polynomial in terms of the elementary symmetric functions of its roots.
\item[\textbf{Fundamental theorem of algebra.}] Every nonconstant polynomial with complex coefficients has at least one complex root.
\item[\textbf{Abel–Ruffini theorem.}] No general formula using radicals exists for polynomial equations of degree five or higher.
\item[\textbf{Bézout's theorem.}] Two projective plane curves of degrees $m$ and $n$ meet in exactly $mn$ points counted with multiplicity.
\item[\textbf{Resultant.}] A polynomial in the coefficients of two polynomials that vanishes when they share a common root.
\item[\textbf{Symmetric function.}] A polynomial in the roots unchanged by any permutation of them.
\end{description}
\section*{7. Sample Questions}
\emph{Exercise reference:} \cite{equations_ref1}. The cited edition is the source for the exercise set; consult its Exercises section for the official statement and numbering.
\begin{enumerate}[leftmargin=*,itemsep=0.9em]
\item \textbf{Exercise 1.} Let $\alpha, \beta, \gamma$ be the roots of $x^3 - 6x^2 + 11x - 6 = 0$. Without finding the roots explicitly, compute $\alpha^2 + \beta^2 + \gamma^2$ using Vieta's formulas. \emph{Solution.}$\alpha + \beta + \gamma = 6$, $\alpha\beta + \beta\gamma + \gamma\alpha = 11$, $\alpha\beta\gamma = 6$. Then $\alpha^2 + \beta^2 + \gamma^2 = (\alpha + \beta + \gamma)^2 - 2(\alpha\beta + \beta\gamma + \gamma\alpha) = 36 - 22 = 14$.

\item \textbf{Exercise 2.} Compute the discriminant of $x^3 - 3x + 1$. Is it zero? \emph{Solution.}$p = -3$, $q = 1$. Discriminant: $\Delta = -4(-3)^3 - 27(1)^2 = 108 - 27 = 81 > 0$. Since $\Delta \neq 0$, the polynomial has three distinct roots. In fact $\Delta > 0$ for a cubic with real coefficients, confirming three distinct real roots.

\item \textbf{Exercise 3.} Compute the discriminant of $x^4 - 2x^3 + x^2 = x^2(x-1)^2$. \emph{Solution.}$\Delta = 256(-2)^3(-1)^4 - 192(-2)^2(1)^2(-1)^2 + \ldots$ Simpler: since $x^2(x-1)^2$ has repeated roots at $x=0$ and $x=1$ (each of multiplicity 2), the discriminant is $\Delta = 0$, as expected.

\item \textbf{Exercise 4.} Verify that $x^5 - x - 1$ is irreducible over $\mathbb{F}_2$ by checking it has no roots and no quadratic factors in $\mathbb{F}_2[x]$. \emph{Solution.} Check roots in $\mathbb{F}_2$: $f(0) = -1 = 1 \neq 0$, $f(1) = 1 - 1 - 1 = -1 = 1 \neq 0$. No linear factors. Check quadratic factors: the only irreducible quadratic over $\mathbb{F}_2$ is $x^2 + x + 1$. Dividing: $x^5 - x - 1 = (x^2 + x + 1)(x^3 + x^2) + (x^2 + x + 1) + 1$, so the remainder is nonzero. Since no irreducible quadratic divides $f$ and $f$ has no roots, $f$ is irreducible over $\mathbb{F}_2$, hence over $\mathbb{Z}$ by Gauss's lemma.

\item \textbf{Exercise 5.} Show that the roots of $x^2 - 2x + 2 = 0$ are complex conjugates, and compute them. \emph{Solution.}$\Delta = 4 - 8 = -4 < 0$, so the roots are complex. By the quadratic formula: $x = \frac{2 \pm \sqrt{-4}}{2} = \frac{2 \pm 2i}{2} = 1 \pm i$. The roots are $1 + i$ and $1 - i$, which are complex conjugates. Their sum is $2$ and product is $(1+i)(1-i) = 1 + 1 = 2$, matching Vieta's formulas.

\end{enumerate}
\section*{8. References}
\begin{thebibliography}{9}
\bibitem{equations_ref1} J. V. Uspensky, \emph{Theory of Equations}, McGraw-Hill, 1948.
\bibitem{equations_ref2} I. Stewart, \emph{Galois Theory}, Chapman and Hall/CRC, 2015.
\bibitem{equations_ref3} D. Cox, J. Little, and D. O'Shea, \emph{Ideals, Varieties, and Algorithms}, Springer-Verlag, 2015.
\end{thebibliography}
